\begin{abstract}
Vehicular networks operating at millimetre-wave and other high-frequency bands
require accurate and timely beam management, yet repeated beam sweeping can
consume substantial radio resources in dynamic road environments. This thesis
develops a Distributed Multimodal Sensing-Assisted Beam Management framework,
denoted by DMSA-BM, for user association and tracking in vehicular networks.
The base-station-centric setting combines
multiview cameras at the serving base station (BS), multiview cameras at distributed
sensing nodes, and a BS-side point cloud derived from integrated sensing and
communication (ISAC). These observations are fused in a bird's-eye-view (BEV)
representation to support vehicle detection, per-target beam prediction, and
selection of a set of \(K\) candidate beams. To bridge
sensing outputs and communication identities, the framework
uses channel state information reference signal (CSI-RS) measurements selected
consistently with their synchronisation signal block (SSB) parents for initial or renewed
target-to-user association and an interacting multiple-model filter to maintain
physical-target tracks. A system-level communication abstraction based on beam
training overhead, regularised zero-forcing precoding, downlink
signal-to-interference-plus-noise ratio (SINR), and
overhead-adjusted effective rate is introduced to evaluate how
sensing, association, and tracking affect communication performance. The thesis
also defines a multimodal simulation and evaluation methodology that connects
road traffic generation, multisensor observations, ray-tracing-based beam
labels, and communication metrics for vehicular networks.
\end{abstract}
